\documentclass[11pt,a4paper]{article}

\usepackage[margin=1in]{geometry}
\usepackage{amsmath}
\usepackage{amssymb}
\usepackage{xcolor}
\usepackage{titlesec}
\usepackage{enumitem}
\usepackage{parskip}
\usepackage{hyperref}
\usepackage{listings}
\usepackage{booktabs}

\definecolor{accentblue}{RGB}{30, 80, 160}
\definecolor{codegray}{RGB}{240,240,240}

\titleformat{\section}{\large\bfseries\color{accentblue}}{\thesection.}{0.5em}{}[\titlerule]
\titleformat{\subsection}{\bfseries}{\thesubsection.}{0.5em}{}

\hypersetup{colorlinks=true, linkcolor=accentblue, urlcolor=accentblue}

\lstset{
  backgroundcolor=\color{codegray},
  basicstyle=\ttfamily\small,
  breaklines=true,
  frame=single,
  framesep=4pt,
  xleftmargin=4pt,
  language=Python,
  showstringspaces=false
}

\title{\color{accentblue}\bfseries QLSTM: Gradients \& Training Mechanics\\[2pt]
\large Reference Sheet for Presentation}
\author{Gautam Gupta --- 2023220}
\date{}

\begin{document}
\maketitle

% =====================================================================
\section{What is Trainable?}

The QLSTM has \textbf{125 total trainable parameters} (with $n_q=2$, $L=1$, hidden$=4$),
divided into three categories. Every one of these is registered with PyTorch via
\texttt{nn.Parameter} or \texttt{nn.Linear}, so they all show up in
\texttt{model.parameters()} and are updated by Adam in a single step.

\subsection{(A) Classical pre-nets --- one per VQC}

\begin{lstlisting}
self.pre_net = nn.Linear(input_size, n_qubits)
\end{lstlisting}

For each gate VQC: $W \in \mathbb{R}^{n_q \times (\text{input}+\text{hidden})}$, $b \in \mathbb{R}^{n_q}$.

With $\text{input}=1, \text{hidden}=4, n_q=2$: weights $5{\times}2 = 10$, bias $2$, $\Rightarrow$ \textbf{12 params per VQC}.

Four VQCs $\Rightarrow$ \textbf{48 classical pre-net params total.}

\subsection{(B) Quantum parameters (the actual circuit angles)}

\begin{lstlisting}
self.q_params = nn.Parameter(torch.randn(n_layers, n_qubits, 3) * 0.1)
\end{lstlisting}

Shape $(L, n_q, 3)$ --- three angles $(\alpha, \beta, \gamma)$ per qubit per layer,
fed into the \texttt{Rot} gate inside \texttt{StronglyEntanglingLayers}:
\[
  \mathrm{Rot}(\alpha, \beta, \gamma) = R_Z(\gamma) \, R_Y(\beta) \, R_Z(\alpha)
\]

With $L=1, n_q=2$: $1 \cdot 2 \cdot 3 = 6$ angles per VQC, four VQCs $\Rightarrow$
\textbf{24 quantum angles total.}

\textit{These are the parameters that suffer from barren plateaus.}

\subsection{(C) Classical post-nets + final readout}

\begin{lstlisting}
self.post_net = nn.Linear(n_qubits, output_size)   # per VQC
self.fc_out   = nn.Linear(hidden_size, output_size) # once at top
\end{lstlisting}

Four post-nets at $4 \cdot (2 \cdot 4 + 4) = 48$ params,\footnote{Pre/post-net output dim equals \texttt{hidden\_size}=4 in our config.} plus \texttt{fc\_out} $= 5$ params.

\subsection{Parameter budget summary}

\[
\underbrace{48}_{\text{4 pre-nets}} \;+\; \underbrace{24}_{\text{quantum angles}} \;+\; \underbrace{48}_{\text{4 post-nets}} \;+\; \underbrace{5}_{\text{fc\_out}} \;=\; \boxed{125 \text{ params}}
\]

\textbf{Only 19\% of the model is quantum.} The rest is classical glue.

% =====================================================================
\section{The Quantum Circuit: Forward Pass Equations}

For one VQC call on input vector $x \in \mathbb{R}^{n_q}$ (already pre-net + tanh'd):

\subsection{Step 1 --- Angle embedding}

\[
  |\psi_0\rangle = \bigotimes_{i=0}^{n_q - 1} R_Y(x_i)\,|0\rangle
  \quad\text{where}\quad
  R_Y(\theta) = \begin{pmatrix}\cos(\theta/2) & -\sin(\theta/2)\\ \sin(\theta/2) & \cos(\theta/2)\end{pmatrix}
\]

\subsection{Step 2 --- Strongly entangling layers}

For each layer $\ell = 1, \dots, L$:
\[
  U_\ell(\theta_\ell) = U_{\text{CNOT-ring}} \cdot \prod_{i=0}^{n_q-1} \mathrm{Rot}(\alpha_{\ell,i}, \beta_{\ell,i}, \gamma_{\ell,i})_i
\]

Full circuit unitary:
\[
  U(\theta) = U_L(\theta_L) \cdots U_2(\theta_2) \, U_1(\theta_1)
\]

Final state:
\[
  |\psi(\theta, x)\rangle = U(\theta)\,|\psi_0(x)\rangle
\]

\subsection{Step 3 --- Measurement}

\[
  z_i(\theta, x) = \langle\psi(\theta, x)|\, Z_i \,|\psi(\theta, x)\rangle \in [-1, 1], \qquad i = 0, \dots, n_q - 1
\]

The output of the VQC is the vector $\mathbf{z}(\theta, x) = (z_0, \dots, z_{n_q-1})^\top$, which then feeds into \texttt{post\_net}.

\paragraph{What is $z_i$, concretely?} After the circuit produces the final state $|\psi(\theta, x)\rangle$, we measure the Pauli-$Z$ observable on each qubit. The number $z_i$ is the \emph{expectation value} of that measurement on qubit $i$, where $Z_i$ means ``Pauli-$Z$ acting on qubit $i$, identity on every other qubit'':
\[
  Z_i \;=\; I \otimes I \otimes \cdots \otimes \underbrace{Z}_{\text{position } i} \otimes \cdots \otimes I
\]

Pauli-$Z$ has eigenvalues $+1$ on $|0\rangle$ and $-1$ on $|1\rangle$, so $z_i$ has a clean probabilistic interpretation:
\[
  z_i \;=\; P(\text{qubit } i = 0) \;-\; P(\text{qubit } i = 1) \;\in\; [-1, +1]
\]
It is a single real number summarising ``is qubit $i$ more likely to be 0 or 1, and by how much.''

\paragraph{What the VQC actually outputs.} A VQC with $n_q$ qubits returns $n_q$ such numbers --- one expectation value per qubit:
\[
  \mathbf{z}(\theta, x) \;=\; \bigl(z_0,\; z_1,\; \dots,\; z_{n_q-1}\bigr) \;\in\; [-1, 1]^{n_q}
\]
In code (\texttt{src/qlstm.py:41}):
\begin{lstlisting}
return [qml.expval(qml.PauliZ(i)) for i in range(n_qubits)]
\end{lstlisting}
That list comprehension is literally producing $[z_0, z_1, \dots, z_{n_q-1}]$.

\paragraph{Why these are what we differentiate.} The vector $\mathbf{z}$ is the \emph{only} thing the quantum circuit hands back to the classical world --- it is the input to \texttt{post\_net}. Everything downstream (post-net $\to$ sigmoid $\to$ cell update $\to$ loss) is classical. So the quantum/classical boundary lives at $\mathbf{z}$, and PennyLane's job is exactly to compute $\partial z_i / \partial \theta_k$ at that boundary. Once those numbers exist, PyTorch's autograd chains them through to $\partial L / \partial \theta_k$.

\paragraph{Concrete tiny example.} Suppose $n_q = 2$ and the final state happens to be the Bell state
\[
  |\psi\rangle = \tfrac{1}{\sqrt{2}}|00\rangle + \tfrac{1}{\sqrt{2}}|11\rangle.
\]
Then
\[
  z_0 = \langle\psi|Z_0|\psi\rangle = \tfrac{1}{2}(+1) + \tfrac{1}{2}(-1) = 0, \qquad
  z_1 = \langle\psi|Z_1|\psi\rangle = 0.
\]
So the VQC outputs $\mathbf{z} = (0, 0)$ --- just two real numbers, which then go into the classical \texttt{post\_net}.

\paragraph{Summary.}
\begin{itemize}
  \item $|\psi\rangle$ = the full $2^{n_q}$-dimensional quantum state, lives inside the simulator.
  \item $z_i$ = a single real number in $[-1, 1]$, the expectation of $Z$ on qubit $i$.
  \item $\mathbf{z}$ = the $n_q$-vector of these numbers, the actual classical output of the VQC.
  \item Gradients $\partial z_i / \partial \theta_k$ are what PennyLane computes; everything else is autograd.
\end{itemize}

\subsection{State evolution --- what does the qubit register actually look like?}

Starting from the all-zero register $|0\rangle^{\otimes n_q}$, the state evolves through three stages:

\textbf{(a) After angle embedding} --- apply $R_Y(x_i)$ on each qubit:
\[
  |\psi_0(x)\rangle \;=\; \bigotimes_{i=0}^{n_q-1} R_Y(x_i)\,|0\rangle
  \;=\; \bigotimes_{i=0}^{n_q-1}\Bigl(\cos\tfrac{x_i}{2}\,|0\rangle + \sin\tfrac{x_i}{2}\,|1\rangle\Bigr)
\]
This is still a \emph{product state}: no entanglement yet, just data loaded into single-qubit amplitudes.

\textbf{(b) After one strongly-entangling layer} --- per-qubit \texttt{Rot} followed by a CNOT ring:
\[
  |\psi_1(\theta_1, x)\rangle \;=\; U_{\text{CNOT-ring}}\,\Bigl[\bigotimes_{i=0}^{n_q-1}\mathrm{Rot}(\alpha_{1,i}, \beta_{1,i}, \gamma_{1,i})\Bigr]\,|\psi_0(x)\rangle
\]
where $U_{\text{CNOT-ring}} = \text{CNOT}_{n_q-1, 0}\,\cdots\,\text{CNOT}_{1, 2}\,\text{CNOT}_{0, 1}$. After this layer the state is generally \emph{entangled} and no longer factorises across qubits.

\textbf{(c) After all $L$ layers} --- the final state going into the measurement step:
\[
  \boxed{\;
  |\psi(\theta, x)\rangle \;=\; \biggl[\prod_{\ell=1}^{L} U_{\text{CNOT-ring}}\,\bigotimes_{i=0}^{n_q-1}\mathrm{Rot}(\alpha_{\ell,i}, \beta_{\ell,i}, \gamma_{\ell,i})\biggr]\;\bigotimes_{i=0}^{n_q-1} R_Y(x_i)\,|0\rangle
  \;}
\]

This is an entangled $n_q$-qubit state living in the full Hilbert space $\mathbb{C}^{2^{n_q}}$, jointly parameterised by the data $x$ and trainable angles $\theta$. It cannot in general be written as a tensor product of single-qubit states --- this is precisely what gives the VQC its expressive power, and also what causes barren plateaus when $n_q$ grows large (the state behaves more and more like a Haar-random vector in $\mathbb{C}^{2^{n_q}}$).

The measurement step in \S2.3 then collapses this $2^{n_q}$-dimensional state down to $n_q$ classical numbers $z_i = \langle\psi|Z_i|\psi\rangle$, which are then handed to the classical \texttt{post\_net}.

% =====================================================================
\section{How Gradients are Computed}

Three different gradient mechanisms run in different parts of the graph.
The chain rule glues them together.

\subsection{Mechanism 1 --- Standard PyTorch autograd (classical parts)}

For \texttt{pre\_net}, \texttt{post\_net}, \texttt{fc\_out}, \texttt{tanh}, \texttt{sigmoid}, the cell-state update --- everything is just tensor ops. PyTorch builds the computation graph during the forward pass and reverse-walks it during \texttt{loss.backward()}.

\subsection{Mechanism 2 --- Quantum gradients via \texttt{diff\_method="backprop"}}

When autograd hits \texttt{self.circuit(x[i], self.q\_params)}, PennyLane takes over. With \texttt{diff\_method="backprop"} on the \texttt{default.qubit} simulator, PennyLane internally represents the circuit as a sequence of unitary matrix multiplications on a state vector, and backpropagates symbolically.

It returns gradients with respect to \emph{both}:
\[
  \frac{\partial z_i}{\partial \theta} \quad\text{and}\quad \frac{\partial z_i}{\partial x}
\]

The second one is what allows the classical pre-net to receive gradient signal.

\paragraph{What is the ``objective'' here, exactly?} The thing being differentiated at this stage is just the measurement output
\[
  z_i(\theta, x) = \langle\psi(\theta, x)|\, Z_i\, |\psi(\theta, x)\rangle
\]
treated as a function of $\theta$ (and of $x$). We are \emph{not} differentiating the loss yet --- the loss-gradient is built up later by the chain rule, using PyTorch autograd to stitch together all the classical pieces. PennyLane's job is just to deliver the partial $\partial z_i / \partial \theta_k$ at the quantum boundary.

\paragraph{How \texttt{backprop} actually computes the gradient.} On a simulator, the quantum state is a complex vector $|\psi\rangle \in \mathbb{C}^{2^{n_q}}$ and every gate is a matrix multiplication $|\psi\rangle \mapsto U_k\,|\psi\rangle$. So the entire circuit is just a sequence of differentiable tensor operations:
\[
  |\psi(\theta, x)\rangle = U_L(\theta_L) \cdots U_2(\theta_2)\,U_1(\theta_1)\,|\psi_0(x)\rangle,
  \qquad
  z_i = \langle\psi|\,Z_i\,|\psi\rangle
\]
PennyLane registers each $U_\ell(\theta_\ell)$ inside PyTorch's autograd graph as a \texttt{torch.Tensor} matrix-multiply with \texttt{requires\_grad=True}. When \texttt{loss.backward()} is called, standard \emph{reverse-mode automatic differentiation} flows through:
\begin{enumerate}
  \item the final $z_i$, giving $\partial z_i / \partial |\psi\rangle$;
  \item the chain $U_L \cdots U_1$, one matrix at a time;
  \item at each $U_\ell(\theta_\ell)$, autograd uses the closed-form derivative (since $U = e^{-i\theta P/2}$ has a known symbolic derivative);
  \item out comes $\partial z_i / \partial \theta_k$ for \emph{all} angles in one reverse sweep.
\end{enumerate}

\paragraph{Why it is exact.} There is no finite-difference approximation anywhere. Every step is an analytic tensor operation, just like differentiating $\sin\theta$ symbolically rather than numerically. The same exactness as classical \texttt{nn.Linear} backprop --- only the operations happen to be complex matrix multiplications instead of real ones.

\paragraph{Why we use it in this project.} On the simulator, \texttt{backprop} is \emph{much} faster than parameter-shift (one reverse sweep vs.\ $2P$ extra forward passes), and gives mathematically identical gradients. Since we are studying trainability and barren plateaus, not benchmarking on real hardware, the cheaper method suffices. See \texttt{src/qlstm.py:34}.

\subsection{Mechanism 3 --- Parameter-shift rule (real hardware)}

If you change \texttt{diff\_method="backprop"} to \texttt{diff\_method="parameter-shift"}, PennyLane computes each quantum gradient via the famous shift identity:
\[
  \boxed{\;\frac{\partial z_i}{\partial \theta_k} \;=\; \frac{1}{2}\Bigl[\, z_i\bigl(\theta + \tfrac{\pi}{2}\,\mathbf{e}_k\bigr) \;-\; z_i\bigl(\theta - \tfrac{\pi}{2}\,\mathbf{e}_k\bigr)\,\Bigr]\;}
\]

\paragraph{What this rule says, in words.} For \emph{each} quantum angle $\theta_k$ that needs a gradient, we run the \emph{entire circuit twice}:
\begin{enumerate}
  \item Once with $\theta_k$ replaced by $\theta_k + \pi/2$, all other angles fixed --- get measurement $z^{(+)}$.
  \item Once with $\theta_k$ replaced by $\theta_k - \pi/2$, all other angles fixed --- get measurement $z^{(-)}$.
\end{enumerate}
Then $\partial z_i / \partial \theta_k = \tfrac{1}{2}(z^{(+)} - z^{(-)})$. That is the gradient. No finite-difference approximation, no Taylor expansion --- it is an algebraic identity, not a numerical estimate.

\paragraph{Why it is exact (one-line proof).} For a gate of the form $U(\theta) = e^{-i\theta P/2}$ where $P$ is a Pauli (so $P^2 = I$):
\[
  U(\theta) = \cos(\theta/2)\,I \;-\; i\sin(\theta/2)\,P
\]
Compute $z(\theta) = \langle\psi(\theta)|\,O\,|\psi(\theta)\rangle$ and take $\partial/\partial\theta$. The result, after using $\sin' = \cos$ and the addition formulas, is \emph{exactly}
\[
  \frac{\partial z}{\partial \theta} = \tfrac{1}{2}\bigl[z(\theta + \pi/2) - z(\theta - \pi/2)\bigr].
\]
The trick is that a $\pm\pi/2$ shift swaps $\sin \leftrightarrow \cos$, so the derivative reduces to a finite combination of shifted evaluations.

\paragraph{Cost.} Two extra full circuit evaluations \emph{per quantum parameter}. With 24 quantum angles in our QLSTM, that is 48 extra forward passes for one backward step. Multiply by sequence length 10 and 4 gates per cell and you understand why running QLSTM on real hardware is brutally slow.

\paragraph{Chain rule still applies.} Parameter-shift only computes $\partial z_i / \partial \theta_k$. To get the loss gradient, the standard chain rule kicks in:
\[
  \frac{\partial L}{\partial \theta_k} \;=\; \underbrace{\frac{\partial L}{\partial z_i}}_{\text{from PyTorch autograd}} \cdot \underbrace{\frac{\partial z_i}{\partial \theta_k}}_{\text{parameter-shift rule}}
\]
PennyLane handles both pieces and hands the final scalar gradient to PyTorch.

\subsection{\texttt{backprop} vs.\ parameter-shift --- side-by-side}

\begin{center}
\begin{tabular}{@{}lll@{}}
\toprule
                            & \texttt{backprop}                          & parameter-shift \\
\midrule
Where it works              & Simulator only                              & Simulator + real hardware \\
Cost / backward             & 1 reverse-mode sweep                        & $2P$ extra circuit runs \\
Memory                      & $O(2^{n_q})$ (state vector stored)          & $O(1)$ (only run circuits) \\
Exactness                   & Exact (symbolic differentiation)            & Exact (algebraic identity) \\
Why exact                   & Direct chain rule on matrix ops             & $\sin \leftrightarrow \cos$ at $\pm\pi/2$ \\
Real-hardware compatible    & No                                           & Yes \\
\bottomrule
\end{tabular}
\end{center}

\textbf{One sentence each:}
\begin{itemize}
  \item \texttt{backprop} is ``I have full access to the simulator's state vector, so I'll just chain-rule through the matrix multiplications.''
  \item parameter-shift is ``I have only a black-box quantum device that runs circuits and measures; I'll exploit the algebraic identity $\sin \leftrightarrow \cos$ to recover exact gradients with two extra runs per parameter.''
\end{itemize}

Both produce the same numerical gradient. We use \texttt{backprop} because we are not on hardware.

% =====================================================================
\section{The Full Chain Rule (End-to-End)}

For a single training sample, sequence length $T$, the gradient of the loss
$L$ with respect to one quantum angle $\theta_k$ inside \texttt{vqc\_forget} is:

\[
  \frac{\partial L}{\partial \theta_k} = \sum_{t=1}^{T}
  \underbrace{\frac{\partial L}{\partial \hat{y}}}_{\text{loss head}}
  \cdot \underbrace{\frac{\partial \hat{y}}{\partial h_T}}_{\text{fc\_out}}
  \cdot \underbrace{\frac{\partial h_T}{\partial h_t}}_{\text{BPTT}}
  \cdot \underbrace{\frac{\partial h_t}{\partial f_t}}_{\text{sigmoid + cell update}}
  \cdot \underbrace{\frac{\partial f_t}{\partial \mathbf{z}_t}}_{\text{post\_net}}
  \cdot \underbrace{\frac{\partial \mathbf{z}_t}{\partial \theta_k}}_{\text{PennyLane}}
\]

\textbf{Reading the chain (right to left):}
\begin{enumerate}
  \item \texttt{loss.backward()} starts at the loss.
  \item Autograd flows through \texttt{fc\_out}, the unrolled time loop (BPTT), the cell-state update, the sigmoid, and the post-net --- all classical.
  \item At the boundary $\partial \mathbf{z}_t / \partial \theta_k$, \textbf{PennyLane handles the quantum part} via \texttt{diff\_method}.
  \item Gradient also flows back through $\partial \mathbf{z}_t / \partial x$ into the pre-net.
\end{enumerate}

\subsection{BPTT --- backprop through time}

Because \texttt{QLSTM.forward} runs a Python loop \texttt{for t in range(seq\_len)} (line 139), the cell is unrolled $T$ times in the computation graph. The gradient $\partial h_T / \partial h_t$ recursively expands as:
\[
  \frac{\partial h_T}{\partial h_t} = \prod_{\tau=t+1}^{T} \frac{\partial h_\tau}{\partial h_{\tau-1}}
\]

With $T = 10$ and 4 gates, each backward pass executes the quantum-gradient call \textbf{$10 \times 4 = 40$ times per sample}.

% =====================================================================
\section{The Barren Plateau Problem (mathematically)}

\subsection{The killer result --- McClean et al.\ 2018}

For a randomly-initialised parameterised quantum circuit with $n$ qubits drawn from a unitary 2-design:
\[
  \mathbb{E}_{\theta}\!\left[\frac{\partial L}{\partial \theta_k}\right] = 0
  \quad\text{and}\quad
  \boxed{\;\mathrm{Var}_{\theta}\!\left[\frac{\partial L}{\partial \theta_k}\right] \in O\!\left(\frac{1}{2^n}\right)\;}
\]

The mean is zero (so gradients are signed noise), and the variance \emph{decays exponentially} with qubit count.

\subsection{What this means for Adam}

Adam's update rule:
\[
  \theta \leftarrow \theta - \eta \cdot \frac{\hat m_t}{\sqrt{\hat v_t} + \varepsilon}
  \qquad
  \hat m_t \approx \mathbb{E}[g_t], \;\; \hat v_t \approx \mathbb{E}[g_t^2]
\]

If $\mathrm{Var}[g] \to 0$, then $\hat v_t \to 0$ and gradients are dominated by numerical noise. Adam moves the quantum angles by negligible amounts $\Rightarrow$ they stay frozen near random init $\Rightarrow$ training fails.

\textbf{Crucially, the classical parameters keep updating} --- they're not subject to BP. So you'd see the loss drop a little (from classical fits) and then stall, with the quantum block never having contributed.

\subsection{Empirical measurement (from \texttt{src/bp\_full\_grid.py})}

For each $(n, L)$ pair:
\begin{enumerate}
  \item Sample $\theta \sim U[0, 2\pi]^P$ and a random input.
  \item Compute one forward + backward pass to get $\nabla_\theta L \in \mathbb{R}^P$.
  \item Repeat 200 times.
  \item Compute per-parameter variance, then mean across parameters:
  \[
    \widehat{\mathrm{Var}}[\nabla L] = \frac{1}{P} \sum_{k=1}^{P}
      \underbrace{\frac{1}{N-1} \sum_{s=1}^{N} \bigl(g_k^{(s)} - \bar g_k\bigr)^2}_{\text{sample variance for parameter } k}
  \]
\end{enumerate}

\textbf{Why mean across parameters and not just $\theta_0$:} if $\theta_0$ controls a $Z$-rotation immediately before a $Z$ measurement, then $[Z, Z] = 0$ gives a structural zero gradient that has nothing to do with BP. Averaging avoids this artefact.

% =====================================================================
\section{The Optimiser Step}

Once gradients exist, the optimiser doesn't care where they came from:
\begin{lstlisting}
optimizer = torch.optim.Adam(model.parameters(), lr=0.02)
loss.backward()      # populates .grad on all 125 Parameters
optimizer.step()     # one Adam update for all of them
\end{lstlisting}

\texttt{model.parameters()} returns all classical \emph{and} quantum leaf tensors indistinguishably. Adam applies the same update rule to every one. There is no special ``quantum optimiser.''

% =====================================================================
\section{Sanity-check Snippet (Live Demo)}

Run this during the talk to show all gradients flow:

\begin{lstlisting}
model = QLSTM(input_size=1, hidden_size=4, output_size=1,
              n_qubits=2, n_layers=1)
loss = model(torch.randn(2, 10, 1)).sum()
loss.backward()

for name, p in model.named_parameters():
    print(f"{name:50s} shape={tuple(p.shape)}  "
          f"grad_norm={p.grad.norm().item():.4f}")
\end{lstlisting}

You will see:
\begin{itemize}
  \item All 125 parameters have non-zero gradients $\Rightarrow$ training works end-to-end.
  \item \texttt{q\_params} gradients are smaller than classical ones $\Rightarrow$ foreshadows BP.
  \item Each VQC's pre-net, q\_params, and post-net appear separately in the list.
\end{itemize}

% =====================================================================
\section{One-Slide Summary}

\begin{center}
\begin{tabular}{@{}lll@{}}
\toprule
\textbf{Component} & \textbf{Gradient by} & \textbf{Count} \\
\midrule
4 pre-nets       & PyTorch autograd                         & 48 params \\
4 VQC angle sets & PennyLane (\texttt{backprop} or shift)   & 24 params \\
4 post-nets      & PyTorch autograd                         & 48 params \\
1 fc\_out        & PyTorch autograd                         & 5 params \\
\midrule
\textbf{Total}   & all glued by chain rule                  & \textbf{125 params} \\
\bottomrule
\end{tabular}
\end{center}

\textbf{The whole hybrid network trains with one \texttt{loss.backward()} + one \texttt{optimizer.step()}.}
PennyLane is the only ``unusual'' piece, and it just provides the gradient
$\partial \mathbf{z} / \partial \theta$ at the quantum boundary. Everything else
is classical deep learning.

% =====================================================================
\section{Two-Stage Least-Squares Mitigation --- Simple Explanation}

\subsection{The setup in one picture}

The QLSTM has this shape:
\[
  \text{input } X \;\xrightarrow{\text{QLSTM cell}}\; h_T \;\xrightarrow{\texttt{fc\_out}}\; \hat y
\]
Everything before \texttt{fc\_out} is non-linear (and contains the quantum stuff). \texttt{fc\_out} is just one linear layer:
\[
  \hat y \;=\; \texttt{fc\_out}(h_T) \;=\; w^\top h_T + b
\]

\textbf{Key observation:} if we \emph{freeze} the QLSTM cell, then training the model is just \emph{linear regression} from $h_T$ to $\hat y$. Linear regression has a closed-form solution. No gradient descent needed.

\subsection{What is $\Phi$?}

$\Phi$ is just the matrix of \textbf{$h_T$ values}, one row per training sample.

For our S\&P 500 setup:
\begin{itemize}
  \item Training samples: $N = 80$
  \item Hidden size: $d = 4$
\end{itemize}

So for each training sample $i$, the QLSTM eats the 10-day price sequence and spits out a 4-dimensional vector $h_T^{(i)} \in \mathbb{R}^4$. Stack all 80 of them:
\[
  \Phi \;=\;
  \begin{pmatrix}
    h_T^{(1)\top} \\
    h_T^{(2)\top} \\
    \vdots \\
    h_T^{(80)\top}
  \end{pmatrix}
  \;\in\; \mathbb{R}^{80 \times 4}
\]

\textbf{In one sentence:} $\Phi$ is what comes out of the QLSTM, before the final linear layer, for every training sample.

\subsection{What is $y$?}

$y$ is the matrix of \textbf{true labels} you want to predict. For S\&P 500, the label is the next-day price (one number per sample):
\[
  y \;=\;
  \begin{pmatrix}
    y^{(1)} \\
    y^{(2)} \\
    \vdots \\
    y^{(80)}
  \end{pmatrix}
  \;\in\; \mathbb{R}^{80 \times 1}
\]

\textbf{In one sentence:} $y$ is just the training labels (true next-day prices), stacked.

\subsection{The closed-form solution}

We want to find weights $w \in \mathbb{R}^{4 \times 1}$ such that
\[
  \Phi w \;\approx\; y.
\]
Schematically:
\[
  \underbrace{\begin{pmatrix} \cdot\;\cdot\;\cdot\;\cdot \\ \cdot\;\cdot\;\cdot\;\cdot \\ \vdots \\ \cdot\;\cdot\;\cdot\;\cdot \end{pmatrix}}_{\Phi:\;80 \times 4}
  \;
  \underbrace{\begin{pmatrix} w_1 \\ w_2 \\ w_3 \\ w_4 \end{pmatrix}}_{w:\;4 \times 1}
  \;\approx\;
  \underbrace{\begin{pmatrix} y^{(1)} \\ y^{(2)} \\ \vdots \\ y^{(80)} \end{pmatrix}}_{y:\;80 \times 1}
\]

This is just linear regression. Solving the ridge version (with $\lambda I$ for stability):
\[
  \boxed{\;w^* \;=\; \bigl(\Phi^\top \Phi + \lambda I\bigr)^{-1}\,\Phi^\top y\;}
\]

That gives the best 4 numbers $(w_1, w_2, w_3, w_4)$. We copy them into \texttt{fc\_out.weight}. Done with Stage 1.

\subsection{Stage 2 --- now train normally}

After Stage 1, we have a great \texttt{fc\_out}. Stage 2 just runs normal Adam training on \emph{all} parameters (quantum + classical), exactly like the random-init baseline. The only difference: we start from a better point.

\subsection{Why it fails on S\&P 500}

Here is the simple version of the problem:

\textbf{Stage 1 fits $w$ to the features that the \emph{starting} (random) quantum parameters happen to produce.}

Call those features $\Phi_0$.

\textbf{Stage 2 changes the quantum parameters.}

Once the quantum parameters move, the features change too. Call them $\Phi_t$. So:
\[
  w^* \;\;\text{was optimal for}\;\; \Phi_0, \qquad \text{but now we have}\;\; \Phi_t \neq \Phi_0.
\]

The carefully-tuned \texttt{fc\_out} is no longer matched to the features it sees. We have basically broken the warm-start. On S\&P 500 (which is non-stationary, so over-fitting is dangerous), this hurts:

\begin{center}
\begin{tabular}{@{}lc@{}}
\toprule
\textbf{Method} & \textbf{Test MAE} \\
\midrule
Random init QLSTM         & 1.242 \\
Identity-block init QLSTM & 1.246 \\
\textbf{Two-stage LS warm-up} & \textbf{2.374} \;(worse!) \\
\bottomrule
\end{tabular}
\end{center}

\subsection{What we did novel --- the fix}

The diagnosis is: \emph{features drift away from where Stage 1 solved.}

The fix is: \textbf{add a penalty that keeps the features from drifting.}

We \textbf{save $\Phi_0$} after Stage 1. Then during Stage 2, at every training step, we compute the current features $\Phi_t$ and add a penalty if they have moved:

\[
  \boxed{\;
  L_{\text{total}}(\theta) \;=\;
  \underbrace{\|\hat y - y\|^2}_{\text{usual MSE}}
  \;+\;
  \lambda_{\text{feat}} \cdot \underbrace{\|\Phi_t - \Phi_0\|^2}_{\text{don't drift!}}
  \;}
\]

In words: ``Train normally, but also pay a price every time the QLSTM features stray from where they were in Stage 1.''

\paragraph{In code} (\texttt{src/mitigations.py:train\_with\_feature\_penalty}):
\begin{lstlisting}
def hook(mod, inp, out):
    captured["features"] = inp[0]   # KEEP gradient (no .detach())

handle = model.fc_out.register_forward_hook(hook)

for _ in range(epochs):
    pred       = model(X)                          # forward pass
    data_loss  = loss_fn(pred, y)
    Phi_t      = captured["features"]              # current features
    feat_pen   = ((Phi_t - Phi_0) ** 2).mean()
    total      = data_loss + feature_lambda * feat_pen
    total.backward()                               # gradient flows
    opt.step()                                     # quantum + classical updated
\end{lstlisting}

The \emph{key} line: in Stage 1's hook we used \texttt{inp[0].detach()} (we just wanted the numbers). In Stage 2 we do \textbf{not} detach --- so the penalty's gradient flows back through PennyLane into the quantum angles. Adam then nudges $\theta_q$ to keep the features stable.

\subsection{Why this is novel}

\begin{itemize}
  \item The original Boabang \& Gyamerah paper (2026) does Stage 2 as plain Adam --- exactly the recipe that broke for us.
  \item None of the 5 papers I cite use a feature-stability penalty. It came directly from looking at \emph{my own} failed experiment and asking ``what is going wrong?''
  \item It is implemented and runnable in \texttt{src/mitigations.py} (functions \texttt{two\_stage\_ls\_novel\_warmup} and \texttt{train\_with\_feature\_penalty}).
\end{itemize}

\paragraph{Three behaviours from one knob.} The hyper-parameter $\lambda_{\text{feat}}$ controls everything:
\begin{itemize}
  \item $\lambda_{\text{feat}} = 0$ $\;\Rightarrow\;$ original (broken) two-stage LS.
  \item $\lambda_{\text{feat}} = \infty$ $\;\Rightarrow\;$ features frozen $\;\Rightarrow\;$ only \texttt{fc\_out} can move (pure linear regression on random quantum features).
  \item $\lambda_{\text{feat}}$ moderate $\;\Rightarrow\;$ keeps the warm-start \emph{and} lets the quantum block adapt slowly.
\end{itemize}

\subsection{Side-by-side}

\begin{center}
\begin{tabular}{@{}p{0.42\textwidth}p{0.42\textwidth}@{}}
\toprule
\textbf{Original two-stage LS} & \textbf{Our novel variant} \\
\midrule
Stage 1: solve $w^* = (\Phi_0^\top\Phi_0 + \lambda I)^{-1}\Phi_0^\top y$ & Same Stage 1, but \emph{also save} $\Phi_0$ \\
Stage 2: plain Adam on MSE & Stage 2: Adam on $\text{MSE} + \lambda_{\text{feat}} \|\Phi_t - \Phi_0\|^2$ \\
Quantum features drift freely & Quantum features pulled back toward Stage-1 \\
Hurts on S\&P 500 (MAE 2.374) & Designed to keep the warm-start gain \\
\bottomrule
\end{tabular}
\end{center}

\paragraph{One-line takeaway.} Stage 1 is just ridge regression: $\Phi$ = QLSTM features for all training samples, $y$ = true labels, $w^* = (\Phi^\top\Phi + \lambda I)^{-1}\Phi^\top y$. Stage 2 of the original paper breaks the Stage 1 fit because the features drift. Our novel fix saves $\Phi_0$ and adds $\lambda_{\text{feat}}\|\Phi_t - \Phi_0\|^2$ to the loss so they cannot drift.

\end{document}
